\documentclass[a4paper,12pt]{article}
\usepackage[utf8]{inputenc}
\usepackage[T1]{fontenc}
\usepackage[ngerman]{babel}
\usepackage{amsmath}
\usepackage{amssymb}
\usepackage{enumitem}
\usepackage{geometry}
\geometry{a4paper, margin=2.5cm}
\usepackage{parskip}

\DeclareMathOperator{\diag}{diag}

\title{Einsendeaufgaben -- Aufgabe 5.4\\[0.5em]
\large Modul 61112: Lineare Algebra}
\author{}
\date{}

\begin{document}
\maketitle

\section*{Aufgabe 5.4}

Sei $A \in M_{nn}(\mathbb{C})$ eine beliebige Matrix der Form
\[
  A = \diag(a_1, a_2, \ldots, a_n)
\]
Wir beobachten, dass das Einsetzen einer diagonalen Matrix in ein Polynom
$p \in \mathbb{C}[X]$ wieder zu einer diagonalen Matrix resultiert.

Wir definieren also $B := \diag(b_1, \ldots, b_n)$ und suchen die Variablen
$b_1, \ldots, b_n \in \mathbb{C}$, sodass
\[
  A = f(B)
\]
gilt.

Wir erkennen, dass für jedes $b_j$ mit $1 \leq j \leq n$ eine Polynomgleichung
aufgestellt werden kann:
\begin{align*}
  \diag(a_1, \ldots, a_n)
    &= f(\diag(b_1, \ldots, b_n)) \\
    &= \sum_{i=0}^{m} c_i \diag(b_1, \ldots, b_n)^i
       && \left( f := \sum_{i=0}^{m} c_i X^i \right) \\
    &= \sum_{i=0}^{m} c_i \diag(b_1^i, \ldots, b_n^i) \\
    &= \diag\left( \sum_{i=0}^{m} c_i b_1^i, \ldots, \sum_{i=0}^{m} c_i b_n^i \right)
\end{align*}

Da $f$ ein nicht konstantes Polynom ist, existieren nach dem Fundamentalsatz der
Algebra Lösungen der Gleichungen
\[
  0 = \left( \sum_{i=0}^{m} c_i b_j^i \right) - a_j
\]
für $b_j$ mit $1 \leq j \leq n$.

Noch zu untersuchen, ob es auch für nicht diagonalisierbare Matrizen gilt. Wir
wählen $f = X^2$, $n = 2$ und $A = \begin{pmatrix} -1 & 1 \\ 1 & -1 \end{pmatrix}$.
Wir setzen $B := \begin{pmatrix} a & b \\ c & d \end{pmatrix}$ in $f$ ein:
\begin{align*}
  f(B) &= \begin{pmatrix} a & b \\ c & d \end{pmatrix}^2 \\
       &= \begin{pmatrix} a^2 + cb & ab + bd \\ ca + cd & bc + d^2 \end{pmatrix} \\
       &= \begin{pmatrix} a^2 + cb & b(a+d) \\ c(a+d) & bc + d^2 \end{pmatrix} \\
       &= \begin{pmatrix} a^2 + 1 & 1 \\ 1 & (1 \pm a)^2 + 1 \end{pmatrix}
\end{align*}
denn wegen $b(a+d) = 1$ und $c(a+d) = 1$ folgt $b = c = a + d = 1$ oder
$b = c = a + d = -1$. Wegen $a^2 + 1 > -1$, existiert also kein
$B \in M_{22}(\mathbb{C})$ mit $f(B) = A$, weshalb die Voraussetzung der
Diagonalisierbarkeit nicht weggelassen werden darf.

\end{document}
